\chapter{Conclusion}
\label{chap:conclusion}

\section{Research Question Revisited}

This project asked whether evaluation of stateful LLM agents could be expressed in a software-testing style, with readable scenarios that combine assertions over replies, tool trajectories, and application state, while preserving enough structured evidence to diagnose failures.

Existing tools generally specialise in output metrics, benchmark environments, or trace observability. The project therefore investigated whether these capabilities could be brought together in a reusable developer-facing library.

\section{Contributions}

The first contribution is a unified behavioural specification. A scenario is a conversation script with checks woven in after each turn. User lines, tool expectations, output rubrics, and state postconditions live together in one file. The same language covers deterministic matchers and selective LLM rubrics where nuance is needed, so reply, trajectory, and state expectations are declared as parts of one episode rather than scattered across datasets, scorers, and transforms.

The second contribution is a framework-neutral execution and diagnostic model. Integrations for LangChain and Pydantic AI normalise traces so scenarios do not depend on one framework's message format. When an assertion fails, the framework attaches a structured counterexample that names which check broke, on which turn, at which field path, with expected and actual summaries and enough transcript context to act. A CLI and local web UI persist runs, compare experiments, and show the same failure card a developer sees in the terminal.

The third contribution is an integrated discovery-to-regression workflow. Repeated runs support stochastic agents. User simulation and mutation fuzzing explore dialogue a human author might not write, using the same oracles throughout. Shrinking and regression extraction turn an exploratory failure into a small, permanent script the suite can rerun. The library does not try to replace every metric-oriented library on its own ground. It offers a single behavioural contract per episode, precise failure reports at the assertion that broke, and a path from discovery to regression within the same scenario model.

\section{Interpretation of the Findings}

The evaluation evidence supports treating the behavioural episode, not the final reply alone, as the correct unit of evaluation. A polite answer can coexist with a skipped authentication step, a tool called in the wrong order, or a store left in the wrong shape. On TelcoSupportBench, many failures appeared in trace and state lanes while the reply still sounded acceptable. Multi-surface scenarios therefore exposed errors that output-only grading would have missed.

Diagnostics are part of evaluation quality, not an afterthought. A test is more useful when it identifies the failed behavioural expectation rather than only producing a score. The expressiveness and diagnostic studies showed that this integration reduced authoring overhead and produced more localised failures than the compared implementations in the evaluated specimens.

Generated exploration complements authored tests rather than replacing them. Manual scenarios provide stable known cases against which new behaviour can be measured. Simulation and fuzzing explore interaction paths and wording a human author might not script, with little extra authoring effort or developer time beyond the base scenario. Regression extraction bridges exploration and maintenance by pinning discovered failures as permanent scenarios under the same oracles. Together, these studies provide evidence that generative testing can complement manually authored scenarios when teams need both repeatability and broader coverage.

\section{Limitations and Future Work}

Chapter~\ref{chap:evaluation} identified limits to construct, external, and rerun validity. Clarity and diagnostic grades relied on LLM judges. Some lanes depended on LLM criteria. Only one reference agent and one telecom domain were studied, and extracted scenarios did not always reproduce the same failure signature on immediate rerun. Those caveats bound how strongly the reported figures can be generalised.

Several directions follow from those limits. Evaluation on additional domains and agent frameworks would test whether the scenario model transfers beyond TelcoSupportBench. Human assessment of authoring clarity and diagnostic usefulness would strengthen construct validity. Signature-stable replay after extraction, and more robust matching of stochastic trajectories, would make discovery findings easier to keep as regressions. Larger suites and team-based studies would also clarify how the workflow scales in everyday development.

\section{Final Remarks}

Agent reliability needs more than leaderboard scores or final-answer grades. It needs behavioural specifications a team can read, run, debug, and extend. \emph{agent-spec-kit} provides a promising basis for further investigation along those lines. It is explicit about what should happen, precise about what went wrong, and designed to turn exploratory failures into tests the suite can keep.
